\section{The complete ordered weak-C2 shape R5 is admitted}
\label{sec:r17-r5}

\begin{theorem}
Normalize $d=1$. Retain the original physical C2 source, its stationary contact, every strict
source gap, all coefficient signs, the positive-infinitesimal total-scale
intersection and every original prior-admission exclusion. Suppose
$B,z,y,r>0$ are infinitesimal and $t=y/B$, $h=H/B$, $l=1-h=L/B$,
and $j=1+zr-l/t$ are positive bounded units. Then the entire original
four-state parent is a finite tensor SOS at every independently chosen
projective testing pair, with both signs, literal zero/infinite angles,
and real or complex states. These are precisely the unchanged R5 source
hypotheses. There is no comparison of source and testing rates.
\end{theorem}

\paragraph{Source entry and all gaps.}
Put $a=z/B$ and $g=G/B=1-a(1-r-H)$. Since $r,H$ are infinitesimal,
$g>0$ implies $a$ bounded and $a_0\in[0,1]$. Write $\rho=a_0$,
$\chi=h_0\in(0,1)$, $\tau=t_0>1-\chi$ and $\ell=1-\chi$,
$\delta=\tau-\ell>0$. The inequality for $\tau$ follows from $j$
being a positive unit. The exact source formulas WE1 give jointly
\[
\begin{gathered}
 (nr)_0={1\over2(1+\rho)},\quad
 A_0={\tau\over2\chi(1+\rho)},\quad
 C_0={\chi+\tau-\rho\ell\over2\ell(1+\rho)},\\
 D_0={1\over2(1+\rho)\delta},\quad
 (\alpha/B^2)_0={\chi+\tau-\rho\ell\over2(1+\rho)}.
\end{gathered}       \tag{R5.1}
\]
All displayed quantities are positive units; $n$ is unbounded and
$\alpha$ infinitesimal. Thus literal $n=0$ is impossible on this shape,
by the exact identity $r(2n+1)\asymp1$, rather than being discarded.
If $\rho<1$, $\mu=m/a^2$ is bounded with residue
$1/[2(1-\rho^2)]$. If $\rho=1$, $g$ is infinitesimal,
$m$ unbounded and $(mg)_0=1/4$. These statements retain arbitrary
independent $r,g,B$ openings. In particular the source is automatically
in the positive-infinitesimal total-scale regime. No nonemptiness of
the residual intersection is asserted.

\paragraph{An exact centered-excess identity for the entire Gram.}
Use the unnormalized version of the established complete WA3--WA8 identity,
with roots
\[
 \xi=X+ruZ,\quad\eta=Y-huW,\quad
 R_1=uW-t^{-1}vZ,\quad R_2=vW-a^{-1}uZ,\quad T=uvW-BZ.
\]
Let $G$ be its five-by-five matrix, so $F/4={\cal R}^TG{\cal R}$.
Both complete normal rows and the complete hard block are included.
Write the original $m$ event as $M=vW-X+uvY-uZ$. Its centered row is
\[
 R_{mc}=M-gvW=-\xi+uv\eta+(1-g)R_2+huT.             \tag{R5.2}
\]
The equality uses $1-g=a(1-r-H)$ and $H=hB$. It vanishes at the
actual contact. Put $L_m=(-1,uv,0,1-g,hu)^T$.
For any positive real $K<m$, there is the exact matrix identity
\[
 G=G\big|_{m=K}+(m-K)L_mL_m^T+(m-K)g\,H_m,          \tag{R5.3}
\]
where the only nonzero entries of $H_m$ are
\[
 (H_m)_{14}=(H_m)_{41}=-1,\quad
 (H_m)_{25}=(H_m)_{52}=v,\quad
 (H_m)_{44}=2-g,\quad
 (H_m)_{45}=(H_m)_{54}=hu.                           \tag{R5.4}
\]
All other actual source coefficients, especially $\alpha/B^2$, remain
fixed in $G|_{m=K}$. This is a formal matrix substitution, not an assertion
that a modified source is physical or satisfies contact. Equivalently,
at the original event level,
\[
 mM^2-(m-K)R_{mc}^2
 =KM^2+2(m-K)g\,vWM-(m-K)g^2v^2W^2.                \tag{R5.5}
\]
The retained square contains every radial and normal jet of the excess.
Equations R5.3--R5.4 prove boundedness of the remainder directly; no
Hessian-subtraction premise is imported.

\paragraph{All infinitesimal or zero testing angles.}
If $\rho\in(0,1)$, apply the pure-source congruence replacing $\xi$ by
$\widehat\xi=\sqrt n\xi$. Every coefficient is bounded: $nr$ is
bounded, $nr^2$ infinitesimal, $a,t$ bounded units, and $A,C,D,m$ and
$\alpha/B^2$ bounded. At $u=v=0$ the full residue consists of a unit
$\widehat\xi$ pivot, a positive $m_0$ pivot on $R_2$, a positive
$(\alpha/B^2)_0$ pivot on $T$, and the $\eta,R_1$ block
\[
 \begin{pmatrix}
 A_0+C_0+D_0+1/2&-(C_0+D_0+1/2)(1-\chi)+A_0\chi\\
 -(C_0+D_0+1/2)(1-\chi)+A_0\chi&
 (C_0+D_0+1/2)(1-\chi)^2+A_0\chi^2
 \end{pmatrix}.                                    \tag{R5.6}
\]
Its determinant is $A_0(C_0+D_0+1/2)>0$.

If $\rho=0$, replace $R_2$ by $aR_2$ as in WE4--WE5.
The exact divisibility $(b_2-\alpha/B^2)/a$ from WE3 holds for arbitrary
$h$ and $r$; its canceled denominator is a unit at $a=r=B=0$ with
$h,l,t,j$ units. Hence every coefficient remains bounded, including
all full normal entries. The $R_2$ pivot becomes $\mu_0=1/2$;
R5.6 and the other pivots remain positive. No $h/a$ bound is used.

If $\rho=1$, retain $(m-K)R_{mc}^2$ using R5.3, for any fixed
positive real $K$. Since $a$ is a unit, all coefficients in the
remainder are bounded after $\widehat\xi=\sqrt n\xi$.
The constant $\widehat\xi,R_2$ coupling is infinitesimal, while
its $R_2$ pivot has residue $K+2(mg)_0=K+1/2$.
All other zero-angle pivots are R5.6, $1$, and $(\alpha/B^2)_0$.
Thus the complete remainder has strict residue in every source case.
The finite separated Gram lift WA10 applies for independent infinitesimal
$u,v$, including zeros and either sign. Adding the retained exact square
proves the complete local parent without any state restriction.

\paragraph{The finite remote determinant when $\rho<1$.}
For bounded tests, normalize $\widehat X=\sqrt nX$ in the original
parent. Its hard residue is $(1+U)(1+V)\widehat X^2$, and the remaining
three-state form on $(W,Y,Z)$ is
\[
\begin{split}
 P={}&m_0(vW+uvY-uZ)^2+C_0(1+V)(uW-Y)^2
   +A_0(1+U)(1+V)Y^2+D_0(uW-Y-vZ)^2\\
 &+\tfrac12\{(uW-Y-vZ)^2+(1+U)(vY-Z)^2
        +(1+V)U Z^2-(1+U)(1+V)Z^2\},                
\end{split}\tag{R5.7}
\]
where here $u,v$ denote the real testing residues and $U=u^2,V=v^2$.
Set
$\mathcal D=8\chi(1-\rho)(1+\rho)^3\ell\delta>0$.
The exact determinant of its symmetric matrix is
\[
 { (1+U)(1+V)\over\mathcal D}
 \left[(\chi+\tau)(\tau U-\rho V)^2
 +(\chi+\tau-\rho\ell)
 \{\tau\delta U^2V+[\chi+(1-\rho)\tau]UV^2\}\right]. \tag{R5.8}
\]
For $U,V>0$ it is strictly positive. For $U>0,V=0$ it is strictly
positive; for $U=0,V>0$ it is strictly positive when $\rho>0$.
The restriction of R5.7 to $Z=0$ is strictly positive in these cases:
its signed core there is
$\tfrac12[(uW-Y)^2+(1+U)VY^2]$, and the $A,C,m$ packets supply
the stated two-state rank. This restriction and R5.8 prove strict
positive definiteness of the entire three-state residue. All source
coefficients are bounded, so WA10 around the real testing residues
proves these complete remote charts.

\paragraph{Finite remote tests with the unbounded m opening.}
Let $\rho=1$ and extract the exact square in R5.5 before normalizing
$\widehat X=\sqrt nX$. Its remainder has bounded coefficients.
At a real testing pair, its three-state residue is
\[
 B_{\rm fin}+K M_0^2+2\kappa vW M_0,\qquad
 M_0=vW+uvY-uZ,\quad\kappa=1/4,                     \tag{R5.9}
\]
where $B_{\rm fin}$ is R5.7 with the $m_0$ packet omitted and
$A_0,C_0,D_0$ evaluated at $\rho=1$.
For $U,V>0$, restrict to $M_0=0$, so $W=uZ/v-uY$.
Its two-state determinant equals $V^{-1}\lim_{\rho\uparrow1}
\det(P)/m_0$. Formula R5.8 gives
\[
 { (1+U)(1+V)\over16\chi\ell\delta V}
 \left[(\chi+\tau)(\tau U-V)^2
 +(\chi+\tau-\ell)\{\tau\delta U^2V+\chi UV^2\}\right]>0.\tag{R5.10}
\]
At $Z=0$ on that kernel, $W=-uY$ and the positive $A,C$ terms,
together with the nonnegative signed-core restriction, give a positive
$Y^2$ coefficient. Hence $B_{\rm fin}|_{M_0=0}$ is positive definite.
When $u=0,v\ne0$, the kernel is $W=0$ and its form is
$(A_0+C_0+1/2)(1+V)Y^2+D_0(Y+vZ)^2\succ0$.
When $v=0,u\ne0$, the kernel is $Z=0$ and its form is
$(C_0+D_0+1/2)(uW-Y)^2+A_0(1+U)Y^2\succ0$.
For each such real test pair, elementary real completion on the
transverse coordinate $M_0$ now gives a sufficiently large positive
real $K$ for which R5.9 is strictly positive definite. This choice is
legal: $m$ is unbounded, so $m-K>0$ for every real $K$; $K$ never
compares a source rate with a testing rate. Use that $K$ in the exact
full-parent identity R5.5 and apply WA10 to the bounded remainder.

\paragraph{The small-a first-axis chart.}
The only finite remote case left is $\rho=0$, $u_0=0$, $v_0\ne0$.
Use the five roots
\[
 S=aW,\quad R=uW,\quad \widehat X=\sqrt nX,\quad Y,\quad Z
\]. Write the original $m$ packet as
$\mu(vS-aX+auvY-auZ)^2$ and replace $uW$ everywhere else by $R$.
This is exactly the complete identity WE8 with unnormalized $Y$ and
$A(1+u^2)(1+v^2)Y^2$ retained. All coefficients are bounded.
The residue has separate pivots $\mu_0v_0^2$ on $S$ and
$1+v_0^2$ on $\widehat X$. Its remaining $(R,Y,Z)$ form is
\[
 C_0(1+V)(R-Y)^2+A_0(1+V)Y^2+D_0(R-Y-vZ)^2
 +\tfrac12\{(R-Y)^2+VY^2-2vRZ\}.                  \tag{R5.11}
\]
Its restriction to $Z=0$ is positive definite and its determinant is
\[
 {V^2(\chi+\tau)^2(1+V)\over8\chi\ell\delta}>0.    \tag{R5.12}
\]
This supplies the required strict full residue at the literal axis too.

\paragraph{Every reciprocal chart and literal infinity.}
For unbounded $u$ put $p=1/u$ and use the polynomial parent $p^2F$.
If $v$ is bounded, the $\widehat X$ residue is $1+V$ and the other
three-state residue is exactly
\[
 [C_0(1+V)+D_0+1/2]W^2+A_0(1+V)Y^2
                         +(m_*+1/2)(vY-Z)^2,         \tag{R5.13}
\]
which is strict for every real $v$. Here $m_*=m_0$ when bounded,
and $m_*=K$ after centered extraction if $g$ is infinitesimal.
The additional centered correction in R5.5 has a factor $p$ and
vanishes in residue. This covers $p=0$ literally.

For unbounded $v$, write $q=1/v$ and use $q^2F$. If $u$ is bounded,
the $\widehat X$ pivot is $1+U$ and the other residue is
\[
 m_*(W+uY)^2+C_0(uW-Y)^2+(A_0+1/2)(1+U)Y^2+D_0Z^2.\tag{R5.14}
\]
For $\rho<1$ this is strict unless $\rho=u_0=0$. For $\rho=1$,
the bounded correction is $2\kappa W(W+uY)$; its restriction to
$W+uY=0$ is strict, and a sufficiently large real $K$ again makes
the entire residue strict. If $\rho=u_0=0$, use $S=aW,R=uW$,
$\widehat X,Y,Z$ as above in the complete WE10 identity. Its residue is
\[
 \widehat X^2+\mu_0S^2+C_0(R-Y)^2+(A_0+1/2)Y^2+D_0Z^2\succ0.
                                                               \tag{R5.15}
\]
Every coefficient of this entire polynomial chart is bounded.

Finally, if both angles are unbounded, use $p=1/u,q=1/v$ and
$p^2q^2F$, with the five roots $\widehat X,W,Y,P=pZ,Q=qZ$.
The exact signed-core part after homogeneity is one half of
\[
 (qW+X-pqY-P)^2+(1+p^2)(Y-Q)^2
                    +(1+q^2)(p^2X^2+2XP-P^2).       \tag{R5.16}
\]
All other original packets are retained in their homogeneous forms;
in particular the $m$ packet is $(pW-pqX+Y-Q)^2$ and its centered
row subtracts $gpW$. After $\widehat X=\sqrt nX$ and any needed
centered extraction, the complete residue is
\[
 \widehat X^2+C_0W^2+A_0Y^2+D_0P^2
                                  +(m_*+1/2)(Y-Q)^2\succ0.       \tag{R5.17}
\]
Its coefficients are bounded; the centered correction again contains
$p$ and vanishes in residue. Both $p=0$ and $q=0$ are literal
polynomial values. The roots $P,Q$ need not be independent to use
this positive Gram.

\paragraph{Join and tensor legality.}
The local chart, finite remote cases, the one small-a finite axis,
each single reciprocal chart and the double reciprocal chart exhaust
all independently chosen testing types. Every use of WA10 whitens a
strict pure-source matrix at real testing residues and splits finitely
many bounded-source times infinitesimal-testing coefficients into
positive pure-source/pure-testing products. There is no mixed inverse
or radical. The roots themselves use only pure-source radicals such
as $\sqrt n$ and pure-testing monomials. Exact positive excess squares
are restored in every chart. Away from literal endpoints, homogeneous
restoration uses only pure-testing squares; at the endpoints the
polynomial value is used. The real identities apply separately to real
and imaginary state parts. All four original state coordinates, both
normal directions, all hard couplings, arbitrary source openings and
the unchanged physicality/total-scale/prior-exclusion intersection
have therefore been retained. This is an all-tests admission of R5.
